% Blueprint: Hatcher
% Created by Numina

\begin{document}

\section{Topology of cell complexes}

These results concern CW complexes, formalized over Mathlib's CW-complex API
(\texttt{Topology.CWComplex}), in which a CW structure is a predicate on a set
$C$ in an ambient space $X$ with explicit cells and characteristic maps. A
\emph{subcomplex} is a subset that is a union of cells whose closures it
contains; it is \emph{finite} when it has finitely many cells.

\subsection{Compactness (Proposition A.1)}

The proof of Proposition~A.1 is decomposed into four steps over Mathlib's
subcomplex API: a finite union of finite subcomplexes is a finite subcomplex
(Lemma~\ref{lem:cw-finite-union}); a finite subcomplex containing the frontier
of a cell can be enlarged to contain the whole closed cell
(Lemma~\ref{lem:cw-insert-cell}); every closed cell lies in a finite subcomplex
(Lemma~\ref{lem:cw-cell-finite-subcomplex}); and a compact subset meets only
finitely many open cells (Lemma~\ref{lem:cw-compact-finite-cells}).

\begin{lemma}[Finite union of finite subcomplexes]
    \label{lem:cw-finite-union}
    \lean{HatcherCW.iUnion_subcomplex_finite}
    \leanfile{AlgebraicTopology/HatcherCW.lean}
    \leanok
    Let $C$ be a CW complex in a Hausdorff space $X$. Given a finite family
    $(E_i)_{i \in s}$ of finite subcomplexes of $C$, there is a finite
    subcomplex $G$ of $C$ with $E_i \subseteq G$ for every $i \in s$.
\end{lemma}
\begin{proof}
    The underlying set of a subcomplex is the union of its open cells and is
    closed, and the cell-index of a union of two subcomplexes is the union of
    their cell-indices. Hence the pointwise union $E \cup F$ of two subcomplexes
    is again a subcomplex, whose cells are the union of those of $E$ and $F$; if
    both are finite, so is the union. Iterating over the finite index set $s$
    gives the required finite subcomplex $G$.
\end{proof}

\begin{lemma}[Enlarging a subcomplex by one cell]
    \label{lem:cw-insert-cell}
    \lean{HatcherCW.subcomplex_insert_cell}
    \leanfile{AlgebraicTopology/HatcherCW.lean}
    \leanok
    Let $C$ be a CW complex in a Hausdorff space $X$, let $E$ be a subcomplex,
    and let $e^n_i$ be an $n$-cell whose frontier $\partial e^n_i$ is contained
    in $E$. Then there is a subcomplex $F$ with $E \subseteq F$ and
    $\overline{e^n_i} \subseteq F$; if $E$ is finite, so is $F$.
\end{lemma}
\begin{proof}
    Take $F$ to be the subcomplex whose cell-index is that of $E$ together with
    the single cell $i$ in dimension $n$, and whose underlying set is
    $E \cup \overline{e^n_i}$. Since $\overline{e^n_i} = e^n_i \cup
    \partial e^n_i$ and $\partial e^n_i \subseteq E$, the closed cells of all
    indexed cells lie in $F$, so this defines a subcomplex. It contains $E$ and
    $\overline{e^n_i}$, and adds a single cell to $E$, so it is finite when $E$
    is.
\end{proof}

\begin{lemma}[Each closed cell lies in a finite subcomplex]
    \label{lem:cw-cell-finite-subcomplex}
    \lean{HatcherCW.closedCell_subset_finite_subcomplex}
    \leanfile{AlgebraicTopology/HatcherCW.lean}
    \uses{lem:cw-finite-union, lem:cw-insert-cell}
    \leanok
    Let $C$ be a CW complex in a Hausdorff space $X$. For every cell $e^n_i$ of
    $C$ there is a finite subcomplex $F$ with $\overline{e^n_i} \subseteq F$.
\end{lemma}
\begin{proof}
    \uses{lem:cw-finite-union, lem:cw-insert-cell}
    By strong induction on $n$. The frontier $\partial e^n_i$ is contained in a
    finite union of closed cells of dimension $< n$. By the induction
    hypothesis each of these closed cells lies in a finite subcomplex, and by
    Lemma~\ref{lem:cw-finite-union} their union lies in a single finite
    subcomplex $E$ containing $\partial e^n_i$. Applying
    Lemma~\ref{lem:cw-insert-cell} to $E$ and the cell $e^n_i$ yields a finite
    subcomplex containing $\overline{e^n_i}$.
\end{proof}

\begin{lemma}[Compact subsets meet finitely many cells]
    \label{lem:cw-compact-finite-cells}
    \lean{HatcherCW.finite_cells_meeting_compact}
    \leanfile{AlgebraicTopology/HatcherCW.lean}
    \leanok
    Let $C$ be a CW complex in a Hausdorff space $X$ and $K \subseteq C$
    compact. Then the set of cells $e^n_i$ of $C$ whose open cell meets $K$ is
    finite.
\end{lemma}
\begin{proof}
    Otherwise choose a point $x_i \in e^n_i \cap K$ in each of infinitely many
    distinct cells met by $K$, and let $S = \{x_i\}$. Each closed cell is
    compact and contains only finitely many of the $x_i$ (one per cell it
    meets, and a closed cell meets finitely many open cells via its frontier),
    so $S$ meets every closed cell in a finite, hence closed, set. By the weak
    topology of the CW complex $S$ is therefore closed in $C$, and the same
    argument applied to each subset of $S$ shows $S$ is discrete. A closed
    discrete subset of the compact set $K$ is finite, contradicting that the
    $x_i$ are distinct.
\end{proof}

\begin{proposition}[Compact subspace lies in a finite subcomplex]
    \label{prop:compact-in-finite-subcomplex}
    \lean{HatcherCW.compact_subset_finite_subcomplex}
    \leanfile{AlgebraicTopology/HatcherCW.lean}
    \uses{lem:cw-finite-union, lem:cw-cell-finite-subcomplex, lem:cw-compact-finite-cells}
    \leanok
    Let $C$ be a CW complex in a Hausdorff space $X$, and let $K \subseteq C$ be
    compact. Then there is a finite subcomplex $E$ of $C$ with $K \subseteq E$.
\end{proposition}
\begin{proof}
    \uses{lem:cw-finite-union, lem:cw-cell-finite-subcomplex,
      lem:cw-compact-finite-cells}
    By Lemma~\ref{lem:cw-compact-finite-cells} the compact set $K$ meets only
    finitely many open cells. Since $K \subseteq C$ is the union of the open
    cells, $K$ is contained in the union of the finitely many closed cells it
    meets. By Lemma~\ref{lem:cw-cell-finite-subcomplex} each of these closed
    cells lies in a finite subcomplex, and by
    Lemma~\ref{lem:cw-finite-union} their union lies in a single finite
    subcomplex $E$, which therefore contains $K$.
\end{proof}

\subsection{Whitehead's characterization (Proposition A.2)}

Proposition A.2 is J.H.C. Whitehead's axiomatic description of CW complexes: a
Hausdorff space $X$ with a family of characteristic maps
$\Phi_\alpha : D^n_\alpha \to X$ carries a CW structure with these as its
characteristic maps iff (i) each $\Phi_\alpha$ is injective on
$\operatorname{int} D^n_\alpha$, the resulting open cells are disjoint and cover
$X$; (ii) each $\Phi_\alpha(\partial D^n_\alpha)$ lies in a finite union of cells
of dimension $< n$; (iii) $X$ has the weak topology with respect to the closed
cells. Mathlib's \texttt{Topology.CWComplex} class is \emph{defined} by exactly
these axioms, together with the characteristic-map data, with one addition: it
also requires continuity of the inverse characteristic map
(\texttt{continuousOn\_symm}). Hatcher derives this --- the ``hence a
homeomorphism'' in condition (i) --- from compactness of $D^n_\alpha$ and
Hausdorffness of $X$, and that derivation is the mathematical content of the
$\Leftarrow$ direction over Mathlib's API.

\begin{lemma}[Inverse characteristic map is continuous]
    \label{lem:whitehead-inverse-homeo}
    \lean{HatcherWhitehead.continuousOn_invFunOn_ball}
    \leanfile{AlgebraicTopology/HatcherWhitehead.lean}
    \leanok
    Let $X$ be Hausdorff and $f : D^n \to X$ continuous on the closed unit ball,
    injective on the open ball, with $f(\operatorname{int} D^n)$ disjoint from
    $f(\partial D^n)$. Then the inverse of $f$ on the open cell
    $f(\operatorname{int} D^n)$ is continuous.
\end{lemma}
\begin{proof}
    Using the weak-topology characterization of continuity, it suffices, for
    each closed $C \subseteq D^n$, to exhibit a closed set whose intersection
    with the open cell equals $f^{-1}$ applied to $C$. Take $f(C \cap D^n)$,
    which is compact (continuous image of a compact set) hence closed in the
    Hausdorff space $X$. Injectivity on the open ball and disjointness of the
    open cell from the sphere image show both this set and the preimage meet the
    open cell in exactly $f(\operatorname{int} D^n \cap C)$.
\end{proof}

\begin{proposition}[Whitehead's characterization of CW complexes]
    \label{prop:whitehead-cw}
    \lean{HatcherWhitehead.mkWhitehead}
    \leanfile{AlgebraicTopology/HatcherWhitehead.lean}
    \uses{lem:whitehead-inverse-homeo}
    \leanok
    Let $X$ be Hausdorff with characteristic maps $\Phi_\alpha$ satisfying
    conditions (i)--(iii) above (with (ii) phrased, as Hatcher states it, in
    terms of open cells of lower dimension). Then the $\Phi_\alpha$ are the
    characteristic maps of a CW structure on $C$.
\end{proposition}
\begin{proof}
    \uses{lem:whitehead-inverse-homeo}
    Conditions (i)--(iii) are precisely the data and axioms of Mathlib's
    \texttt{CWComplex} class except for the inverse-continuity axiom
    \texttt{continuousOn\_symm}. That axiom is supplied by
    Lemma~\ref{lem:whitehead-inverse-homeo}: for each cell, condition (ii) (open
    form) together with disjointness of the open cells (i) shows the open cell is
    disjoint from its sphere image, so the lemma gives continuity of the inverse
    characteristic map. The closed-cell form of (ii) required by the class
    follows from the open-cell form since open cells are contained in closed
    cells.
\end{proof}

\subsection{Normality (Proposition A.3)}

Proposition A.3 of Hatcher's appendix asserts that every CW complex is normal,
and in particular Hausdorff. Over Mathlib's \texttt{Topology.CWComplex} API the
ambient space $X$ is already assumed Hausdorff, so the mathematical content is
the normality of the subspace $C$: any two disjoint closed subsets of $C$ have
disjoint open neighborhoods. Hatcher's proof builds these neighborhoods by an
inductive construction over the skeleta: given a subset $A$ and a choice of
$\varepsilon_\alpha > 0$ for each cell, one defines an open set
$N_\varepsilon(A) \supseteq A$ whose preimage under each characteristic map
$\Phi_\alpha : D^n_\alpha \to X$ is an open $\varepsilon$-neighborhood of
$\Phi_\alpha^{-1}(A) \cup \partial D^n_\alpha$ together with a collar
$(1-\varepsilon, 1] \times \Phi_\alpha^{-1}(N_\varepsilon^{n-1}(A))$ in spherical
coordinates. For disjoint closed $A$ and $B$ the cell-wise radii can be chosen
small enough that $N_\varepsilon(A)$ and $N_\varepsilon(B)$ are disjoint, which
gives normality.

\begin{lemma}[Separating neighborhoods of disjoint closed sets]
    \label{lem:cw-separating-nbhd}
    \lean{HatcherNormal.exists_disjoint_nhds}
    \leanfile{AlgebraicTopology/HatcherNormal.lean}
    Let $C$ be a CW complex in a Hausdorff space $X$. For any two disjoint
    subsets $A, B$ of $C$ that are closed in $C$, there exist disjoint open
    subsets $U, V$ of $C$ with $A \subseteq U$ and $B \subseteq V$.
\end{lemma}
\begin{proof}
    Build the open neighborhoods $N_\varepsilon(A)$ and $N_\varepsilon(B)$ by
    induction over the skeleta as above, starting from the empty neighborhood on
    the $0$-skeleton and extending across each cell by adjoining an open
    $\varepsilon$-neighborhood of the preimage together with a spherical collar.
    Since $A$ and $B$ are closed and disjoint, on each closed cell the preimages
    of $A$ and $B$ are disjoint compact sets, so a positive radius
    $\varepsilon_\alpha$ can be chosen making the cell-wise neighborhoods
    disjoint; the resulting unions $N_\varepsilon(A)$ and $N_\varepsilon(B)$ are
    open, contain $A$ and $B$, and are disjoint.
\end{proof}

\begin{proposition}[CW complexes are normal]
    \label{prop:cw-normal}
    \lean{HatcherNormal.normalSpace}
    \leanfile{AlgebraicTopology/HatcherNormal.lean}
    \uses{lem:cw-separating-nbhd}
    Let $C$ be a CW complex in a Hausdorff space $X$. Then the subspace $C$ is a
    normal space (and, being a subspace of a Hausdorff space, Hausdorff).
\end{proposition}
\begin{proof}
    \uses{lem:cw-separating-nbhd}
    Points are closed in $C$ since they pull back to closed sets under every
    characteristic map, and $C$ is Hausdorff as a subspace of the Hausdorff space
    $X$. Normality is exactly
    Lemma~\ref{lem:cw-separating-nbhd}: disjoint closed subsets of $C$ are
    separated by the disjoint open neighborhoods $N_\varepsilon(A)$ and
    $N_\varepsilon(B)$.
\end{proof}

\section{The Compact-Open Topology}

This section formalizes the core point-set results about the compact-open
topology from the Appendix (\emph{Topology of Cell Complexes}) of Hatcher's
\emph{Algebraic Topology}, specifically Propositions A.13, A.14, and A.16.
Throughout, $C(Y, X)$ denotes the space of continuous maps $Y \to X$ equipped
with the compact-open topology.

\subsection{Uniform convergence (Proposition A.13)}

When the domain is compact and the codomain is a metric space, the compact-open
topology is the topology of uniform convergence, induced by the supremum metric.
This is recorded in Mathlib as the metric space structure on $C(X, Y)$, whose
distance is the pointwise supremum.

\begin{lemma}[Pointwise bound on distance]
    \label{lem:dist-le}
    \lean{Hatcher.dist_le}
    \leanfile{AlgebraicTopology/Hatcher.lean}
    \leanok
    Let $X$ be a compact space and $Y$ a metric space. For $f, g \in C(X, Y)$
    and a real number $C \ge 0$, we have $d(f, g) \le C$ if and only if
    $d(f(x), g(x)) \le C$ for every $x \in X$.
\end{lemma}
\begin{proof}
    The metric on $C(X, Y)$ is pulled back along the uniform embedding of
    $C(X, Y)$ into the bounded continuous functions $X \to Y$ (which is an
    isometry since $X$ is compact). The corresponding statement for bounded
    continuous functions is exactly that the distance is bounded by $C$ iff
    every pointwise distance is. Unfolding this equivalence through the
    embedding gives the result.
\end{proof}

\begin{proposition}[Compact-open topology is uniform convergence]
    \label{prop:dist-eq-isup}
    \lean{Hatcher.dist_eq_iSup}
    \leanfile{AlgebraicTopology/Hatcher.lean}
    \uses{lem:dist-le}
    \leanok
    Let $X$ be a nonempty compact space and $Y$ a metric space. For
    $f, g \in C(X, Y)$,
    \[
        d(f, g) = \sup_{x \in X} d(f(x), g(x)).
    \]
    Thus the compact-open topology on $C(X, Y)$ coincides with the metric
    topology of uniform convergence.
\end{proposition}
\begin{proof}
    \uses{lem:dist-le}
    By Lemma~\ref{lem:dist-le} the number $d(f, g)$ is an upper bound for the
    pointwise distances $x \mapsto d(f(x), g(x))$, and it is the least such
    upper bound: any $C$ with $C \ge d(f(x), g(x))$ for all $x$ satisfies
    $C \ge d(f, g)$. Hence $d(f, g)$ is the supremum of the pointwise
    distances. The supremum is well defined because $X$ is compact, so the
    continuous function $x \mapsto d(f(x), g(x))$ is bounded.
\end{proof}

\subsection{Evaluation and the exponential law (Proposition A.14)}

\begin{theorem}[Continuity of evaluation]
    \label{thm:eval-continuous}
    \lean{Hatcher.eval_continuous}
    \leanfile{AlgebraicTopology/Hatcher.lean}
    \leanok
    Let $Y$ be a locally compact space and $X$ any space. The evaluation map
    \[
        e : C(Y, X) \times Y \to X, \qquad e(f, y) = f(y),
    \]
    is continuous.
\end{theorem}
\begin{proof}
    Fix $(f, y)$ and an open neighborhood $U$ of $f(y)$. Since $f$ is
    continuous and $Y$ is locally compact, there is a compact neighborhood $K$
    of $y$ with $f(K) \subseteq U$, i.e. $f \in M(K, U)$. Then
    $M(K, U) \times K$ is a neighborhood of $(f, y)$ mapped into $U$ by $e$, so
    $e$ is continuous at $(f, y)$.
\end{proof}

\begin{lemma}[Uncurrying a continuous map is continuous]
    \label{lem:uncurry-cont}
    \lean{Hatcher.uncurry_continuous}
    \leanfile{AlgebraicTopology/Hatcher.lean}
    \uses{thm:eval-continuous}
    \leanok
    Let $Y$ be locally compact and let $F : Z \to C(Y, X)$ be continuous. Then
    the uncurried map $\widehat{F} : Z \times Y \to X$, $\widehat{F}(z, y) =
    F(z)(y)$, is continuous.
\end{lemma}
\begin{proof}
    \uses{thm:eval-continuous}
    The map $\widehat{F}$ is the composition
    $Z \times Y \xrightarrow{F \times \mathrm{id}} C(Y, X) \times Y
    \xrightarrow{e} X$ of the continuous map $F \times \mathrm{id}$ with the
    evaluation map $e$, which is continuous by
    Theorem~\ref{thm:eval-continuous} since $Y$ is locally compact.
\end{proof}

\begin{lemma}[Currying a continuous map is continuous]
    \label{lem:curry-cont}
    \lean{Hatcher.curry_continuous}
    \leanfile{AlgebraicTopology/Hatcher.lean}
    \leanok
    Let $g : Z \times Y \to X$ be continuous. Then the adjoint map
    $\check{g} : Z \to C(Y, X)$, $\check{g}(z)(y) = g(z, y)$, is continuous.
    No local compactness hypothesis is needed.
\end{lemma}
\begin{proof}
    It suffices to check that $\check{g}^{-1}(M(K, U))$ is open for each
    subbasic open set $M(K, U)$ of $C(Y, X)$. For $z$ in this preimage,
    $g^{-1}(U)$ is an open neighborhood of the compact set $\{z\} \times K$, so
    there are open $V \ni z$ and $W \supseteq K$ with $V \times W \subseteq
    g^{-1}(U)$; then $V$ is a neighborhood of $z$ contained in the preimage.
\end{proof}

\begin{proposition}[Exponential law]
    \label{prop:exponential-law}
    \lean{Hatcher.exponential_law}
    \leanfile{AlgebraicTopology/Hatcher.lean}
    \uses{lem:uncurry-cont, lem:curry-cont}
    \leanok
    Let $Y$ be locally compact. A map $g : Z \times Y \to X$ is continuous if
    and only if its adjoint $\check{g} : Z \to C(Y, X)$,
    $\check{g}(z)(y) = g(z, y)$, is continuous.
\end{proposition}
\begin{proof}
    \uses{lem:uncurry-cont, lem:curry-cont}
    If $\check{g}$ is continuous, then $g = \widehat{\check{g}}$ is continuous
    by Lemma~\ref{lem:uncurry-cont} (using local compactness of $Y$).
    Conversely, if $g$ is continuous then $\check{g}$ is continuous by
    Lemma~\ref{lem:curry-cont}.
\end{proof}

\subsection{Currying as a homeomorphism (Proposition A.16)}

\begin{lemma}[Currying is continuous]
    \label{lem:continuous-curry}
    \lean{Hatcher.continuous_curry}
    \leanfile{AlgebraicTopology/Hatcher.lean}
    \uses{lem:curry-cont, thm:eval-continuous}
    \leanok
    Let $Y$ and $Z$ be locally compact. The currying map
    \[
        \mathrm{curry} : C(Y \times Z, X) \to C(Y, C(Z, X))
    \]
    is continuous.
\end{lemma}
\begin{proof}
    \uses{lem:curry-cont, thm:eval-continuous}
    By Lemma~\ref{lem:curry-cont} applied twice, it suffices to show that the
    fully uncurried map $C(Y \times Z, X) \times Y \times Z \to X$ is
    continuous. After reassociating the product, this is an iterated evaluation
    map, continuous by Theorem~\ref{thm:eval-continuous} since $Y$ and $Z$ are
    locally compact.
\end{proof}

\begin{lemma}[Uncurrying is continuous]
    \label{lem:continuous-uncurry}
    \lean{Hatcher.continuous_uncurry}
    \leanfile{AlgebraicTopology/Hatcher.lean}
    \uses{lem:curry-cont, thm:eval-continuous}
    \leanok
    Let $Y$ and $Z$ be locally compact. The uncurrying map
    \[
        \mathrm{uncurry} : C(Y, C(Z, X)) \to C(Y \times Z, X)
    \]
    is continuous.
\end{lemma}
\begin{proof}
    \uses{lem:curry-cont, thm:eval-continuous}
    By Lemma~\ref{lem:curry-cont}, it suffices to show the uncurried map
    $C(Y, C(Z, X)) \times (Y \times Z) \to X$ is continuous. Reassociating, this
    is a composition of evaluation maps, each continuous by
    Theorem~\ref{thm:eval-continuous} since $Y$ and $Z$ are locally compact.
\end{proof}

\begin{theorem}[Currying homeomorphism]
    \label{thm:curry-homeo}
    \lean{Hatcher.curryHomeo}
    \leanfile{AlgebraicTopology/Hatcher.lean}
    \uses{lem:continuous-curry, lem:continuous-uncurry}
    \leanok
    Let $Y$ and $Z$ be locally compact spaces and $X$ any space. The currying
    map
    \[
        \mathrm{curry} : C(Y \times Z, X) \xrightarrow{\ \cong\ } C(Y, C(Z, X))
    \]
    is a homeomorphism, with inverse $\mathrm{uncurry}$.
\end{theorem}
\begin{proof}
    \uses{lem:continuous-curry, lem:continuous-uncurry}
    The maps $\mathrm{curry}$ and $\mathrm{uncurry}$ are mutually inverse
    bijections, as one checks by evaluating both composites pointwise. Both are
    continuous by Lemma~\ref{lem:continuous-curry} and
    Lemma~\ref{lem:continuous-uncurry}, so $\mathrm{curry}$ is a homeomorphism.
\end{proof}

\end{document}
